\section{总结与延伸}

两段讲座从不同方向解释 Transformer-like computation。SDM 路线从高维几何出发：distributed addresses 的 intersection count 近似指数相似度，进而得到 softmax retrieval；cerebellar circuit 提供候选 read/write implementation。TEM 路线从认知需求出发：分离 reusable structure 与 changing content，用 fast associative memory 绑定，再把 memory read 改写成 modern Hopfield attention。

\subsection{两条路线的统一比较}

下面的表不是把 cerebellum 与 hippocampus 合并成一个脑模型，而是比较它们分别解释 attention 的哪一部分。第一讲更关注 weighting kernel 与存储几何，第二讲更关注 positional/structural state、rapid binding 与 compositional inference。二者共同强调 key/content separation、distributed memory 与 uncertainty about biological implementation。

\begin{longtable}{@{}L{0.18\textwidth}L{0.24\textwidth}L{0.25\textwidth}L{0.21\textwidth}@{}}
\toprule
路线 & 主要问题 & Transformer 对应 & 最强证据 \\
\midrule
SDM geometry & noisy query 怎样按相似度召回记忆 & softmax attention weights & 高维交集近似与数值实验 \\
Cerebellar mapping & 神经电路怎样分离激活、写入与读出 & keys/values 与 pooled retrieval & 解剖可实现性和候选实验 \\
Cognitive map / TEM & 怎样复用关系结构并快速绑定新内容 & recurrent position state + memory & 行为、模型与神经表征比较 \\
Modern Hopfield TEM & associative read 怎样扩容并接入 context & query/key/value attention & 公式 correspondence 与实现收益 \\
\bottomrule
\end{longtable}

\begin{importantbox}{课程真正的“neuroscience-inspired”方法}
不是复制某个脑区名字，而是先找计算问题，再提出具有明确 inductive bias 的模型，最后在数学、任务行为、内部表征、神经数据和干预实验之间建立可证伪联系。Brain analogy 只有能生成新预测时才比隐喻更有价值。
\end{importantbox}

\subsection{对 Transformer 研究的六条启示}

\begin{enumerate}
  \item \textbf{解释 softmax 的几何条件。} 检查 norm、temperature、dimension 与 approximation region，而不是把 softmax 当神秘默认值。
  \item \textbf{区分短期 pattern memory 与长期 distributed memory。} Context precision 和 parameter capacity 面临不同干扰与可追溯性。
  \item \textbf{把 key 与 value 分离写清楚。} 匹配线索和返回内容不同，是 attention、SDM 和 hippocampal binding 的共同机制。
  \item \textbf{让位置表示反映任务结构。} Sequence index 只是最简单结构；graph transition、action 与 relational state 可提供更强 positional dynamics。
  \item \textbf{为快速绑定保留 fast state。} 新环境内容不应总靠慢速参数更新，memory/fast weights 可承担 episode-specific knowledge。
  \item \textbf{把 biological plausibility 变成实验。} 说明要记录哪些 cell、何时写入、怎样干预、什么结果会否定 mapping。
\end{enumerate}

\begin{warningbox}{本讲最容易被夸大的三句话}
“Attention approximates SDM”不等于 attention 与 SDM 在所有参数下完全相同；“小脑可实现 SDM”不等于小脑被证实运行 Transformer；“TEM 约等于 Transformer”不等于 hippocampus 是语言模型。每句话都必须附上比较层级、假设和反例空间。
\end{warningbox}

\subsection{拓展阅读}

\begin{itemize}
  \item Bricken and Pehlevan, \href{https://proceedings.neurips.cc/paper/2021/hash/8171ac2c5544a5cb54ac0f38bf477af4-Abstract.html}{\emph{Attention Approximates Sparse Distributed Memory}（NeurIPS）}
  \item Whittington et al., \href{https://www.sciencedirect.com/science/article/pii/S009286742031388X}{\emph{The Tolman--Eichenbaum Machine}（Cell）}
  \item Whittington et al., \href{https://arxiv.org/abs/2112.04035}{\emph{Relating Transformers to Models and Neural Representations of the Hippocampal Formation}（arXiv）}
  \item Bricken et al., \href{https://openreview.net/forum?id=JknGeelZJpHP}{\emph{Sparse Distributed Memory is a Continual Learner}（OpenReview）}
  \item Behrens et al., \href{https://www.nature.com/articles/s41593-022-01153-y}{\emph{How to Build a Cognitive Map}（Nature Neuroscience）}
\end{itemize}

阅读时可以固定同一套检查表：计算对象是什么、表示如何因子化、memory 如何写入/读取、哪些等价是 exact/approximate、哪些结果来自模型、哪些来自 neural data、什么实验能否定 mapping。这样才能把 neuroscience 转化为模型设计与科学预测，而不是给已有架构贴脑区标签。

\end{document}
